\chapter*{How to use this book}
\addcontentsline{toc}{chapter}{How to use this book}
\markboth{How to use this book}{How to use this book}

\section*{Reading paths}
Chapters build on one another, but not in a single line. Figure~\ref{fig:readingmap}
shows the eight parts, the chapters inside each, and three suggested emphases.
Chapters marked \textsf{C} form the core path that every reader should follow;
\textsf{T} marks chapters for readers who want the formal theory of molecular
computation; \textsf{L} marks chapters for readers heading towards laboratory
design and molecular engineering.

\TikzPlate{frontmatter/figures/reading-map}{A reading map. Each row is one part
of the book; each card is a chapter, with its number, a short title and badges
for the core path (C), theory emphasis (T) and laboratory and engineering
emphasis (L). Within a part, chapters are best read left to right. A ``needs''
note on a card names the key prerequisites from earlier parts; for example,
strand displacement (Chapter 19) needs thermodynamics and kinetics
(Chapters 6--7).}{fig:readingmap}

\begin{description}
\item[If you come from biology or chemistry.] Read Chapter~1, skim
Chapters~5 and 8--10 (much will be familiar, but notice how each mechanism is
turned into an operation with a precise input and output), and read the
\emph{Computing bridge} boxes. Appendix~\ref{app:math} reviews logarithms,
probability, rates and differential equations, and Chapter~3 introduces
complexity from first principles.
\item[If you come from computer science, mathematics or physics.] Read
Chapter~1 and Appendix~\ref{app:biology} before Part~II, and read the
\emph{Biology bridge} boxes. Do not skip Chapters~5--8: the strand
orientation, thermodynamic and kinetic ideas they introduce are what make the
later algorithms physically meaningful.
\item[If you want the research frontier quickly.] Read Chapters~1, 5, 6,
19, 21, 22 and 25, and then the Epilogue.
\end{description}

\section*{Course plans}
A one-semester (14-week) introductory course can follow the core path:
\begin{center}
\small\sffamily
\begin{tabular}{@{}lll@{}}
\toprule
Weeks & Chapters & Theme\\
\midrule
1--2 & 1, 2, 4 & Computing with molecules; Adleman's experiment; resources\\
3--4 & 5, 6, 7 & Strands, hybridisation thermodynamics, kinetics and noise\\
5--6 & 8, 9, (10) & Enzymes, PCR amplification, readout\\
7 & 11, 12 & Sequence design; generate--filter--verify algorithms\\
8--9 & 19, 20 & Strand displacement; reaction networks as programs\\
10--11 & 21, 22 & DNA circuits and neural networks; self-assembly and origami\\
12 & 25 & DNA data storage and coding\\
13 & 30, 31 & From gene regulation to molecular programming in cells\\
14 & Epilogue & Frontiers, open problems, student projects\\
\bottomrule
\end{tabular}
\end{center}
A second semester, or a graduate course, can take either the theory emphasis
(Chapters 3, 13--18, 20, 22, 27, 29 and 32) or the engineering emphasis
(Chapters 10, 11, 21--26, 28 and 31).

\section*{Features of each chapter}
Each chapter opens with its learning goals and the ideas it relies on, and
closes with a summary of key ideas, exercises and worked solutions. Inside the
chapter, a small number of recurring boxes signal their purpose:
\begin{TryThis}
asks you to commit to a prediction before the text derives the answer.
\end{TryThis}
\begin{Pitfall}
names a misconception, a failure mode or a negative control that a careful
experiment or proof must exclude.
\end{Pitfall}
\begin{BioBridge}
supplies the biological or chemical background a computing reader needs at
that point, without assuming a biology course.
\end{BioBridge}
\begin{CSBridge}
supplies the mathematical or computational background a biology reader needs
at that point, without assuming a programming course.
\end{CSBridge}
\begin{Frontier}
summarises current research, stating what was demonstrated and the scope of
that evidence.
\end{Frontier}

Exercises progress from reading a figure and predicting an outcome, through hand
calculation and tracing an algorithm, to implementation and open design
problems. Worked solutions follow each exercise set; open problems carry an
explicit description of what a good answer must contain.

\section*{Visual conventions}\label{sec:conventions}
The illustrations use one visual language throughout
(Figure~\ref{fig:legend}). Colour is always reinforced by shape, line style or
a label, so every figure can be read in greyscale.

\TikzPlate{frontmatter/figures/visual-legend}{The book's visual language.
Strand identity is shown by colour and a half-arrowhead (harpoon) at the
3$'$ end; base pairs are grey rungs; covalent bonds are continuous lines and
hydrogen bonds dashed. Arrow shape encodes meaning: a chemical reaction, a
reversible association, a physical transfer of material, a symbolic
(computational) mapping, or an observation that supports an inference.
Molecule classes have fixed colours: DNA in blues and violets, RNA in coral,
proteins and enzymes in gold, computational objects in teal and measurements
in violet. Drawings are schematic and not to scale unless the caption says
otherwise.}{fig:legend}

\section*{Notation}
\begin{center}
\small
\begin{tabular}{@{}p{0.27\linewidth}p{0.66\linewidth}@{}}
\toprule
Symbol & Meaning\\
\midrule
$5'$, $3'$ & The two chemically distinct ends of a nucleic-acid strand. Sequences
are written $5'\to3'$ unless an alignment says otherwise.\\
$a$, $a^{*}$ & A named sequence domain and its Watson--Crick complement.\\
$\mathrm{rc}(s)$ & Reverse complement of sequence $s$.\\
$N(P)$ & Multiplicity: number of molecular copies of candidate $P$ in a population.\\
$[X]$ & Concentration of species $X$, usually in mol\,L$^{-1}$ (M), nM or \textmu M.\\
$N_{\mathrm A}$ & Avogadro constant, $6.02214076\times10^{23}\,\mathrm{mol}^{-1}$ (exact).\\
$\Delta G^{\circ}$, $\Delta H^{\circ}$, $\Delta S^{\circ}$ & Standard free energy,
enthalpy and entropy change of a reaction.\\
$K$, $K_d$ & Equilibrium (association) constant; dissociation constant $K_d=1/K$.\\
$T_m$ & Melting temperature: where half of the duplex-forming strands are paired.\\
$k$, $k_f$, $k_r$ & Rate constants; units depend on reaction order (e.g.\ M$^{-1}$s$^{-1}$).\\
$a_j(x)$ & Propensity: the event rate of reaction $j$ in a discrete-count state $x$.\\
$\Pr(\cdot)$, $\mathbb E[\cdot]$ & Probability and expectation.\\
$O(\cdot)$ & Asymptotic upper bound on resource growth (Chapter~3).\\
\bottomrule
\end{tabular}
\end{center}
